% system_overview.tex -- Methodology 3.1: rig, twin and roles (rewritten 12 Oct 2026 to the thesis plan of
% 11 Oct). The safety paragraph and its label sec:safety moved to authority.tex (3.5); the nodes figure is replaced
% by fig:owners there. The twin's description moved here from simulation.tex with its label sec:simulation. T4 is the
% 17 Oct: plates clockwise from +x, angle = 30 x plate (Wesley 14 Oct; docs/rig/plate_numbering.md); the 7 Oct
% registry rows number plates anticlockwise (old p is (12 - p) mod 12 here).

\subsection{Rig, twin and roles}
\label{sec:overview}

\subsubsection{The rig}
\label{sec:hardware}

All control runs on one ground computer under ROS~2 in a 6\,$\times$\,6\,$\times$\,4\,m motion-capture cage.
Motion capture is the only sensing: it gives the pose of each UAV, of the ring and of the newcomer's magnet tip, and
nothing on the ring is instrumented. Each UAV's flight controller runs Betaflight in rate mode and receives a throttle
and three body-rate commands over its own ExpressLRS radio link at 100\,Hz; the arm switch and the magnet share the
link. A centralised planner plans the ring and every rod at 10\,Hz (Section~\ref{sec:controller}), a tracking
controller per UAV turns its share into radio commands at 50\,Hz (Section~\ref{sec:tracker}), and a fleet manager
sequences ARM, TAKEOFF, the trajectory, LAND and the reconfiguration events.

\paragraph{UAVs, rods and ring.} Four UAVs on six-cell batteries, each 0.55\,kg as flown with its rod and magnet. Each carries a rigid rod of 0.55\,m and 0.075\,kg with a ball joint at both
ends; the upper joint sits 4\,cm below the UAV's centre. The ring is 0.86\,kg, 0.51\,m across and 40\,mm thick, with
twelve magnet plates every 30$^\circ$ on its top face at a radius of 0.225\,m. Each rod ends in an electromagnet
that catches on a plate only on contact, holds about 9--10\,N straight off the face and less under a pull across
it, and releases under a jerk rather than under a steady load; [the hold against the angle of pull, to be measured
with a spring scale at 0, 30, 45 and 60$^\circ$ off the plate's axis]. The magnets are held on from
start-up so that a quiet radio link cannot drop the ring, and a magnet is switched off by a channel on its UAV's
own link: that is how a release is commanded, and a release the planner is not told of is made the same way. The
newcomer carries the same rod on a tether that goes slack on the rig once it is free of the ring; in the twin
every link, that tether included, is rigid.

\paragraph{Layouts and plate numbering.} A layout is named by the plates its UAVs stand on, counted clockwise
from plate 0 on the ring's $+x$ axis seen from above, so that plate $k$ sits $30k$ degrees clockwise from $+x$
(Figure~\ref{fig:geometry}); every plate number and angle in this thesis follows that rule. (The collaborator's
code counts anticlockwise: his plate 3 is plate 9 here.) The carriers stand on plates 3, 7 and 11, and the fourth
UAV, the one that leaves and rejoins, on plate 9; its departure leaves an even triangle
(Section~\ref{sec:reconfig}).

% fig_geometry.tex -- M3: payload and rod geometry (TikZ). Dimensions from the rig (30 Sep 2026):
% plates at r 0.225 m on a 500 mm ring, 12 plates every 30 deg, rod 0.55 m, pivot 4 cm below the UAV centre.
% 17 Oct: plates numbered clockwise from +x seen from above (angle = 30 x plate, Wesley 14 Oct); the side view at
% the planned elevation of 60 deg (since 7 Oct). Drawn to scale (1 m = 9 cm).
\begin{figure}[h]
\centering
\begin{tikzpicture}[x=9cm, y=9cm, >=Stealth, font=\small]
  % ---------------- (a) top view ----------------
  \begin{scope}
    \fill[gray!12] (0,0) circle (0.25);
    \fill[white] (0,0) circle (0.19);
    \draw[gray!60, thin] (0,0) circle (0.25) (0,0) circle (0.19);
    % the even triangle left after plate 9 leaves
    \draw[densely dashed, gray!70] (30:0.225) -- (150:0.225) -- (270:0.225) -- cycle;
    % plates, numbered clockwise from +x
    \foreach \k in {0,...,11} {
      \pgfmathsetmacro{\a}{-30*\k}
      \node[draw, fill=white, minimum size=0.11cm, inner sep=0pt, rotate=\a] at (\a:0.225) {};
      \node[font=\scriptsize, text=gray!70!black] at (\a:0.29) {\k};
    }
    % the four carrying plates: 3, 7, 9 and 11
    \foreach \k in {3,7,9,11} { \fill[black] (-30*\k:0.225) circle (0.012); }
    \draw[thick, black] (-270:0.225) circle (0.022);
    % centre, axes, rho
    \fill (0,0) circle (0.008);
    \draw[->, thin, gray] (0,0) -- (0.09,0) node[below, font=\scriptsize, text=black] {$x_L$};
    \draw[->, thin, gray] (0,0) -- (0,0.09) node[left, font=\scriptsize, text=black] {$y_L$};
    \draw[->, thick] (0,0) -- (30:0.213) node[pos=0.62, above left, inner sep=1pt] {$\boldsymbol{\rho}_i$};
    \draw[->, gray] (-12:0.335) arc (-12:-48:0.335);
    \node[font=\scriptsize, text=gray!70!black, anchor=west] at (-30:0.345) {clockwise};
    \node[font=\scriptsize, align=left, anchor=north] at (0,-0.33) {$\bullet$ carrying plates 3, 7, 9, 11\quad $\bigcirc$ plate 9 leaves};
    \node[anchor=north west] at (-0.36,0.36) {(a)};
  \end{scope}
  % ---------------- (b) side view of one rod (A plate, P pivot, C UAV centre), elevation 60 deg ----------------
  \begin{scope}[shift={(0.72,-0.33)}]
    \fill[gray!12] (-0.06,0) rectangle (0.0,0.03);
    \draw[gray!60, thin] (-0.06,0) rectangle (0.0,0.03);
    \fill[black] (-0.012,0.03) rectangle (0.012,0.038);
    \coordinate (A) at (0,0.034);
    \coordinate (P) at (0.275,0.5103);
    \coordinate (C) at (0.275,0.5503);
    \draw[very thick] (A) -- (P);
    \fill (A) circle (0.007) (P) circle (0.007);
    \draw[thick] (0.185,0.5503) -- (0.365,0.5503);
    \draw[gray] (0.185,0.5563) ellipse (0.035 and 0.006) (0.365,0.5563) ellipse (0.035 and 0.006);
    \fill (C) circle (0.006);
    \draw[->, thick] (A) -- (0.08,0.1726) node[pos=0.55, above left, inner sep=1pt] {$\mathbf{q}_i$};
    \draw[gray] (0.16,0.034) arc (0:60:0.16);
    \draw[gray, thin] (A) -- (0.22,0.034);
    \node[font=\scriptsize] at (0.20,0.095) {60$^\circ$};
    \node[font=\scriptsize, anchor=north west] at (0.16,0.30) {$l_i = 0.55$\,m};
    \draw[gray, thin, densely dotted] (P) -- (0.455,0.5103) (C) -- (0.455,0.5503);
    \draw[<->, thin] (0.455,0.5103) -- (0.455,0.5503) node[midway, right, font=\scriptsize] {$\boldsymbol{\delta}$: 4\,cm};
    \node[font=\scriptsize, anchor=north] at (-0.03,-0.005) {ring};
    \node[font=\scriptsize, anchor=south] at (0.275,0.575) {UAV $i$};
    \node[anchor=north west] at (-0.08,0.69) {(b)};
  \end{scope}
\end{tikzpicture}
\caption[Payload and rod geometry]{Payload and rod geometry, to scale. (a) Top view of the 0.86\,kg ring
with its twelve magnet plates at a radius of 0.225\,m, numbered clockwise from the $x_L$ axis seen from above, so
that plate $k$ sits $30k$ degrees clockwise from $x_L$. The four-UAV layout uses plates 3, 7, 9 and 11; when the
UAV on plate 9 leaves, the three on plates 3, 7 and 11 form an even triangle (dashed), under which the ring hangs
level. (b) Side view of one rod at the planned elevation of 60$^\circ$: the rod direction $\mathbf{q}_i$ points
from the plate to the UAV, and the rod pivots $\boldsymbol{\delta}$ = 4\,cm below the UAV's centre.}
\label{fig:geometry}
\end{figure}


\pendingfig{fig:rig_photo}{The rig: four UAVs on rods, the ring with its twelve magnet plates, the motion-capture
markers and the newcomer's magnet tip, with a to-scale top and side view.}{a photograph of the rig, on the next rig day}

\subsubsection{The rig twin}
\label{sec:simulation}

Every controller change is flown in simulation before it reaches the rig, with the same launch files and controller
code. The simulator is Gazebo \cite{koenig2004gazebo}, bridged to ROS~2 \cite{macenski2022ros2}; the worlds are
generated from the rig's measured geometry and masses, and each flight starts with the ring on the floor. A flight
controller emulator maps the throttle to thrust through the rig's identified affine map, with each UAV's own offset
and a sagging battery, and closes a rate loop on the simulated gyro; a motion-capture emulator publishes the poses at
120\,Hz. The twin runs at 0.3--0.45 of real time under simulated time, and runs that fall further behind are voided.

\paragraph{What the twin does not model.} The twin leaves out six things, each named where it matters. (i) Loop latency: the rig's
command path carries about 50\,ms through the radio and 30\,ms through the flight controller, and the twin has only
the latency its own bridges model; in an exploratory check, 15 and 30\,ms of command delay added in the twin left the carry at 0.3\,m/s
as it was, and 45\,ms tipped a carrier at takeoff. (ii) Solver
timing: in the twin the planner always finishes before its next tick, because simulated time waits for it; on the
rig it has 60\,ms. (iii) Motion-capture noise: the twin's poses are exact. (iv) Rotor wake on the ring and on the
other UAVs. (v) The slack tether: every link is rigid in the twin. (vi) The magnet's holding limit: a twin magnet
never lets go unless commanded. Until the evening of 8 October the rig's body rates were rounded to 0.1\,rad/s on their way to the
trackers; rig flights are labelled before or after that fix where it matters.

\subsubsection{Roles}
\label{sec:roles}

Two controllers fly in this work, the same in the twin and on the rig. This work's planner and trackers carry the
ring: the floor start, the carry, the detection of a release, the resize, the confirmation of an attach from motion
capture, the switch of the newcomer's radio, the fold-in, the take-back and the fault policy. The collaborator's
controller flies a free UAV: the pickup of an object, the transit, the drop into the ring, the descent onto a plate to
contact, and the lift that the confirmation reads. Where this work's code runs inside the collaborator's stack it is
listed as this work's, [made with the collaborator's written agreement; the wording to be agreed with him and the
supervisor]; his thesis is cited as in preparation.
